\section{Introduction}
\label{sec:intro}

We touch a mouse and a keyboard all day. Sometimes touching them is a problem:
a doctor wearing gloves, a cook with dirty hands, a teacher standing far from
the laptop during a lecture, a shared screen in a public place, or a person
who finds small devices hard to hold. Almost every laptop already has a
webcam, so a mouse and keyboard that work with \emph{only the webcam} would
cost nothing and help all these people.

The idea is not new. Many open-source ``virtual mouse'' programs find the
hand with MediaPipe~\cite{zhang2020mediapipe}, check which fingers are up with
a few fixed rules, and move the cursor with the index fingertip. When we tried
this approach, it worked in a demo but was hard to use for real work. We saw
four problems:
\begin{enumerate}
\item \textbf{Rules break easily.} A slightly tilted hand, a hand far from the
camera, or simply a different person makes the fixed thresholds fail.
\item \textbf{The cursor shakes.} The detected hand points move a few pixels
in every frame even when the hand is still, so small buttons are hard to hit.
\item \textbf{The cursor jumps when you click.} To click, you pinch the thumb
and index finger. This moves the fingertip, and the cursor goes with it, just
when you need it to stay still.
\item \textbf{False clicks.} When the hand changes from one gesture to another,
the detector is wrong for a frame or two, and each wrong ``pinch'' frame
becomes a click.
\end{enumerate}

Our project is about solving these four problems, and about \emph{measuring}
that each fix really helps. We keep MediaPipe for finding the hand, and
everything after it is our own work:

\begin{itemize}
\item \textbf{A learned gesture classifier on cleaned hand points
(\cref{sec:features,sec:mlp}).} Before classification we remove the
effect of hand position, size, tilt and left/right hand. A small neural
network (an MLP with about 20\,k weights) then learns 7 gestures. We test it
on a person it has never seen.
\item \textbf{A stable cursor point (\cref{sec:cursor}).} The cursor follows the
index-finger \emph{knuckle}, which hardly moves during a pinch.
\item \textbf{A smart smoothing filter (\cref{sec:cursor}).} The One-Euro
filter~\cite{casiez2012oneeuro} smooths strongly when the hand is slow and
weakly when it is fast, so we get a steady cursor without delay.
\item \textbf{Voting over time (\cref{sec:voting}).} A gesture counts only
when it is seen in most of the last five frames with good confidence.
\item \textbf{Extras that make it usable (\cref{sec:actions}).} An air keyboard
with a ``hover-lock'' that types the right key, and a one-second
personalization for a new user.
\end{itemize}

The rest of the report follows the same story. \Cref{sec:related} looks at
what others have done. \Cref{sec:method} explains our method step by step,
and why we chose each step. \Cref{sec:exp} tests every step with our own
recorded data and live sessions. \Cref{sec:conclusion} sums up what we learned.
